\documentclass[tikz,border=10pt]{standalone}
\usepackage{tikz}
\usetikzlibrary{shapes, arrows.meta, positioning, fit, backgrounds, calc, shadows, patterns}
\usepackage{amsmath}
\usepackage{fontspec}
% Note: You might need to change the font if compiling locally without these fonts installed.
% \setmainfont{Arial} 

\begin{document}

% Define Colors matching the template
\definecolor{col1_bg}{RGB}{235, 245, 251} % Light Blue
\definecolor{col1_header}{RGB}{174, 214, 241} % Darker Blue
\definecolor{col2_bg}{RGB}{253, 242, 233} % Light Orange
\definecolor{col2_header}{RGB}{245, 194, 160} % Darker Orange
\definecolor{col3_bg}{RGB}{244, 236, 247} % Light Purple
\definecolor{col3_header}{RGB}{210, 180, 222} % Darker Purple

\definecolor{box_fill}{RGB}{255, 255, 255}
\definecolor{text_dark}{RGB}{50, 50, 50}

\begin{tikzpicture}[
    node distance=1cm,
    font=\sffamily,
    % Styles
    header_node/.style={
        minimum width=6cm, 
        minimum height=1.2cm, 
        text=black, 
        font=\bfseries\large, 
        align=center,
        yshift=-0.2cm
    },
    panel_box/.style={
        draw=none,
        fill=box_fill,
        rounded corners=5pt,
        drop shadow={opacity=0.15},
        inner sep=0.3cm
    },
    arrow_style/.style={
        ->, 
        >={Stealth[length=3mm]}, 
        line width=1.5pt, 
        color=black!70
    }
]

% ==========================================================================================
% COLUMN I: Input & Masking Layer (Blue)
% ==========================================================================================

% Background
\node[minimum width=7cm, minimum height=14cm, fill=col1_bg, rounded corners=10pt] (col1) at (0,0) {};
% Header
\node[header_node, fill=col1_header, rounded corners=8pt, anchor=north] (header1) at ($(col1.north)+(0,-0.2)$) {I. 多模态输入与掩码层\\Input \& Masking Layer};

% Content 1: Raw Data
\node[panel_box, minimum width=6cm, minimum height=4cm, anchor=north] (box1_1) at ($(header1.south)+(0,-0.5)$) {};
\node[anchor=north, font=\bfseries] at (box1_1.north) {原始多模态图数据 (Raw Data)};

% Drawing Graph inside box1_1
\begin{scope}[shift={($(box1_1.center)+(-1.5, 0.5)$)}, scale=0.6]
    \foreach \i in {1,...,5} \node[circle, fill=blue!50, draw=white, inner sep=2pt] (n\i) at (\i*72:1.5) {};
    \draw[gray] (n1)--(n2)--(n3)--(n4)--(n5)--(n1);
    \draw[gray] (n1)--(n3) (n2)--(n4);
    \node[below=0.8cm of n4, font=\scriptsize] {Graph $G$};
\end{scope}

% Drawing Text/Image icons
\node[draw, fill=white, minimum width=1cm, minimum height=1.2cm, drop shadow] (text_icon) at ($(box1_1.center)+(1.5, 1)$) {\tiny Text};
\node[draw, fill=white, minimum width=1.2cm, minimum height=1cm, drop shadow] (img_icon) at ($(box1_1.center)+(1.5, -0.5)$) {\tiny Image};

% Content 2: Learnable Masks
\node[panel_box, minimum width=6cm, minimum height=5cm, anchor=north] (box1_2) at ($(box1_1.south)+(0,-1.0)$) {};
\node[anchor=north, font=\bfseries] at (box1_2.north) {可学习掩码 (Learnable Masks)};

% Mask Visualization
\node[matrix, anchor=center] (mask_vis) at (box1_2.center) {
    \node[font=\small] {$M_G \sim \mathcal{U}(0,1)$}; \\
    \node[font=\small] {$M_T \sim \mathcal{U}(0,1)$}; \\
    \node[font=\small] {$M_I \sim \mathcal{U}(0,1)$}; \\
};
\node[below=0.5cm of mask_vis, align=center, font=\small] {Masking Operation:\\ $G' = G \odot \sigma(M_G)$};

% Arrow between boxes
\draw[arrow_style] (box1_1.south) -- (box1_2.north);


% ==========================================================================================
% COLUMN II: Explanation Optimization Layer (Orange)
% ==========================================================================================

% Background
\node[minimum width=7.5cm, minimum height=14cm, fill=col2_bg, rounded corners=10pt, right=0.5cm of col1] (col2) {};
% Header
\node[header_node, fill=col2_header, rounded corners=8pt, anchor=north, minimum width=6.5cm] (header2) at ($(col2.north)+(0,-0.2)$) {II. 解释优化机制层\\Optimization Mechanism Layer};

% Sub-module 1: Consistency (UMA)
\node[panel_box, minimum width=6.5cm, minimum height=3.2cm, anchor=north] (box2_1) at ($(header2.south)+(0,-0.4)$) {};
\node[anchor=north west, font=\bfseries] at (box2_1.north west) {1. 统一模态一致性 (UMA)};

% Visual for UMA
\begin{scope}[shift={($(box2_1.center)+(0,-0.2)$)}]
    \node[circle, fill=orange!30, draw] (hG) at (0, 0.8) {$h_G$};
    \node[circle, fill=blue!20, draw] (hT) at (-1.5, -0.5) {$h_T$};
    \node[circle, fill=green!20, draw] (hI) at (1.5, -0.5) {$h_I$};
    \draw[<->, dashed, thick, red] (hG) -- (hT);
    \draw[<->, dashed, thick, red] (hG) -- (hI);
    \node[font=\scriptsize, align=center] at (0, -1.2) {$\mathcal{L}_{consist} = \sum (1 - \text{CosSim}(h_G, h_m))$};
\end{scope}


% Sub-module 2: Interaction Detection (ID)
\node[panel_box, minimum width=6.5cm, minimum height=3.2cm, anchor=north] (box2_2) at ($(box2_1.south)+(0,-0.4)$) {};
\node[anchor=north west, font=\bfseries] at (box2_2.north west) {2. 交互检测 (Interaction)};

% Visual for ID
\begin{scope}[shift={($(box2_2.center)+(0,-0.3)$)}]
    \node[draw, fill=gray!10, minimum width=1.5cm] (joint) at (-1.5, 0) {$\Phi(G,T,I)$};
    \node[draw, fill=gray!10, minimum width=1.5cm] (sum) at (1.5, 0) {$\sum \Phi(m)$};
    \node[font=\Large] at (0,0) {$\neq$};
    \node[font=\scriptsize, align=center] at (0, -0.8) {$\Delta_{inter} = |\text{Joint} - \text{Sum}|$};
\end{scope}


% Sub-module 3: Domain Regularization (DKR)
\node[panel_box, minimum width=6.5cm, minimum height=3.2cm, anchor=north] (box2_3) at ($(box2_2.south)+(0,-0.4)$) {};
\node[anchor=north west, font=\bfseries] at (box2_3.north west) {3. 领域知识正则化 (DKR)};

% Visual for DKR
\begin{scope}[shift={($(box2_3.center)+(0,-0.2)$)}]
    % Invalid molecule
    \node[circle, fill=red!20, draw, inner sep=1pt] (c1) at (-2, 0) {C};
    \foreach \a in {0, 72, ..., 288} \draw (c1) -- +(\a:0.5);
    \node[below=0.6cm of c1, font=\tiny, red] {Invalid (Deg>4)};
    
    \draw[->, thick] (-1, 0) -- (1, 0) node[midway, above, font=\tiny] {Constraint};
    
    % Valid molecule
    \node[circle, fill=green!20, draw, inner sep=1pt] (c2) at (2, 0) {C};
    \foreach \a in {0, 90, 180, 270} \draw (c2) -- +(\a:0.5);
    \node[below=0.6cm of c2, font=\tiny, green!50!black] {Valid};
\end{scope}


% ==========================================================================================
% COLUMN III: Output Layer (Purple)
% ==========================================================================================

% Background
\node[minimum width=7cm, minimum height=14cm, fill=col3_bg, rounded corners=10pt, right=0.5cm of col2] (col3) {};
% Header
\node[header_node, fill=col3_header, rounded corners=8pt, anchor=north] (header3) at ($(col3.north)+(0,-0.2)$) {III. 解释生成与输出层\\Output \& Explanation Layer};

% Content 1: Loss Aggregation
\node[panel_box, minimum width=6cm, minimum height=3cm, anchor=north] (box3_1) at ($(header3.south)+(0,-0.5)$) {};
\node[anchor=north, font=\bfseries] at (box3_1.north) {多目标损失聚合 (Loss Aggregation)};
\node[align=center, font=\small] at (box3_1.center) {$\mathcal{L}_{total} = \mathcal{L}_{pred} + \lambda_1 ||M||_1$ \\ $+ \lambda_2 \mathcal{L}_{ID} + \lambda_3 \mathcal{L}_{domain} + \lambda_4 \mathcal{L}_{consist}$};

% Content 2: Output Graph
\node[panel_box, minimum width=6cm, minimum height=6cm, anchor=north] (box3_2) at ($(box3_1.south)+(0,-1.0)$) {};
\node[anchor=north, font=\bfseries] at (box3_2.north) {最终解释子图 (Explanatory Subgraph)};

% Drawing Final Subgraph
\begin{scope}[shift={($(box3_2.center)+(0, 0.5)$)}, scale=0.8]
    % Background faint graph
    \foreach \i in {1,...,6} \node[circle, fill=gray!20, draw=none, inner sep=1.5pt] (bg\i) at (\i*60:1.5) {};
    \draw[gray!20] (bg1)--(bg2)--(bg3)--(bg4)--(bg5)--(bg6)--(bg1);
    
    % Highlighted Motif (e.g., NO2)
    \node[circle, fill=red!50, draw=black, inner sep=2pt] (h1) at (0,0) {N};
    \node[circle, fill=red!50, draw=black, inner sep=2pt] (h2) at (1,0.5) {O};
    \node[circle, fill=red!50, draw=black, inner sep=2pt] (h3) at (1,-0.5) {O};
    \draw[thick, black] (h1)--(h2);
    \draw[thick, double, black] (h1)--(h3);
    
    \node[below=1.5cm of h1, font=\small, align=center] {Identified Motif\\(e.g., Mutagenic Group)};
\end{scope}

% Arrows connecting columns
\draw[arrow_style] (col1.east) -- (col2.west);
\draw[arrow_style] (col2.east) -- (col3.west);

\end{tikzpicture}
\end{document}
